% ECONOMICS_PROBLEM: {"id": "H63-344", "title": "State-contingent sovereign borrowing", "jel_code": "H63", "source_id": "OP-0258"}
% !TeX program = xelatex
\documentclass[11pt,a4paper]{article}
\usepackage[margin=23mm]{geometry}
\usepackage{fontspec}
\setmainfont{Latin Modern Roman}
\usepackage{amsmath,amssymb,mathtools}
\usepackage{xcolor,enumitem,booktabs,longtable,array,xurl,hyperref}
\definecolor{muted}{HTML}{53636C}
\hypersetup{colorlinks=true,urlcolor=blue,linkcolor=blue}
\setlength{\parindent}{0pt}
\setlength{\parskip}{5pt}
\newcommand{\R}{\mathbb R}
\newcommand{\E}{\mathbb E}
\newcommand{\N}{\mathbb N}
\newcommand{\ind}{\mathbf 1}
\newcommand{\Var}{\operatorname{Var}}
\newcommand{\Cov}{\operatorname{Cov}}
\newcommand{\supp}{\operatorname{supp}}
\DeclareMathOperator*{\argmin}{arg\,min}
\DeclareMathOperator*{\argmax}{arg\,max}
\newcommand{\entrymeta}[3]{{\sffamily\small\color{muted}\textbf{\texttt{#1}}\quad|\quad #2\par #3\par}}
\newcommand{\Problemref}[1]{\texttt{#1}}
\newcommand{\JELref}[2]{\texttt{#2} (JEL \texttt{#1})}
\begin{document}
\section*{State-contingent sovereign borrowing}
\textbf{Problem ID:} \texttt{H63-344}\quad \textbf{JEL:} \texttt{H63}\par
% BEGIN ECONOMICS STATEMENT
\label{problem:OP-0258}
{\sffamily\small\color{muted}Original subject group: Fiscal policy and sovereign debt\par}
\label{definitions:problem:30007539}
\textbf{Audit classification:} Explicit published open question. Limited indexation is expressly a puzzle with existing explanations; this contract design is editorial. Verification concerns the cited version, not first priority or proof that this exact editorial specification remains unsolved on October 8, 2026.
\subsection*{Definitions and precise research target}
\paragraph{Editorial mathematical specification.} Fix a sovereign with stochastic true output \(y_t\), observable output report \(\tilde y_t\), and an ability to influence reporting at a real cost. A debt contract promises repayment \(g(\tilde y_t)\), possibly subject to a maximum payment. Investors have a declared stochastic discount factor, ambiguity set for output and reporting, and a liquidity cost depending on contract standardization. The sovereign’s preferences and default technology are given. Let \(\mathcal G\) be a finite-dimensional, publicly stated family of repayment schedules.
Determine whether a schedule \(g\in\mathcal G\) yields a welfare-improving equilibrium relative to noncontingent debt while satisfying lender break-even and truthful-reporting incentives. In particular, require
\[
E[m_t g(\tilde y_t)]-\text{liquidity and ambiguity costs}\ge \text{issue price},
\]
where \(m_t\) is the investor discount factor, together with a specified sovereign incentive-compatibility condition for every feasible false report. Quantify how much insurance survives pricing premia, manipulation, and reduced market liquidity; explain conditions under which noncontingent issuance is an equilibrium despite gross risk-sharing benefits. Assess contract standardization or independent verification as feasible remedies. The problem concerns this explicitly modeled adoption puzzle, not whether state-contingent debt can provide insurance in an ideal frictionless model, a benefit already established.

For a one-period contract let the issue price be \(Q(g)\), actual repayment after default be \(\widehat g(\tilde y,y)\le g(\tilde y)\), and the lender's actual pricing functional be \(\inf_{P\in\mathcal P}E_P[m\widehat g]-c_{\rm liq}(g)\). Competitive zero-profit pricing equates that functional to \(Q(g)\), rather than treating the issue price as a free parameter. The report incentive requires \(u(y-g(y))\ge u(y-g(r))-k(y,r)\) for every feasible false report \(r\), in the no-default one-period benchmark with \(k\) in utility units; the full default model replaces this by continuation values and actual repayments. Contracts must keep consumption positive or specify treatment at zero. A welfare comparison holds the required resources raised fixed, or optimizes borrowing in both regimes; comparing different issue proceeds can confound insurance with financing volume. Require a strictly positive lender discount factor with finite discounted repayments and define all liquidity costs in the issue-price numeraire.
\subsection*{Variants and assumption changes}
The following are \emph{editorial variants}, not additional published open conjectures.
\begin{enumerate}
\item \emph{Verified GDP indexation.} An independent statistical authority observes output without manipulation and the lender has one known law. Compare optimal smooth repayments with noncontingent debt, keeping funding proceeds and default technology identical.
\item \emph{Ambiguous indexation and standardization.} Allow an explicitly bounded ambiguity set and a liquidity cost depending on contract features. Compare continuous schedules with threshold warrants and assess whether standardization lowers costs enough to reverse the ranking.
\end{enumerate}
\subsection*{References and source audit}
\begin{itemize}
\item Leonardo Martinez; Francisco Roch; Francisco Roldán; Jeromin Zettelmeyer. \emph{Sovereign Debt}. 2022. Locator: Section10, State-contingent debt, printed pp.25--26. Primary text checked. \url{https://fqroldan.github.io/resources/MRRZ.pdf}.
\end{itemize}
Limited indexation is expressly a puzzle with existing explanations; this contract design is editorial.
\begin{enumerate}\raggedright
\item\hypertarget{definitions:source:30007539:1}{} Leonardo Martinez; Francisco Roch; Francisco Roldán; Jeromin Zettelmeyer. 2022. \emph{Sovereign Debt}. Section10, State-contingent debt, printed pp.25--26.
\url{https://fqroldan.github.io/resources/MRRZ.pdf}.
DOI: \url{https://doi.org/10.5089/9798400213250.001}.
\end{enumerate}

\subsection*{Common conventions: Source mathematical conventions}

These conventions apply to each entry unless explicitly replaced.

\textbf{Models and identification.} A primitive model \(m\) specifies all probability laws, feasible choices, information, equilibrium and measurement rules used by its target. The admissible class \(\mathcal M\) is fixed independently of outcomes. If \(P_O(m)\) is its observable probability law and \(\tau(m)\) the target, its identified set at observable law \(P\) is
\[
 \mathcal I_\tau(P)=\{\tau(m):m\in\mathcal M,\ P_O(m)=P\}.
\]
Point identification means that this set is a singleton. An empty set signals incompatibility of data and assumptions. Coordinate bounds need not describe the joint set. Bounds are sharp only if every retained target is attainable or arbitrarily approachable by compatible models. Finite bounds require explicit restrictions on unobserved quantities; unrestricted confounding, hidden wealth, extrapolation or arbitrary equilibrium selection can yield uninformative sets.

\textbf{Causal interventions.} A potential outcome \(Y_i(p)\) is the outcome under a complete policy regime \(p\), including timing, financing, information and market spillovers. The target population and units are fixed before treatment. With interference, use the complete assignment vector or a justified exposure mapping; individual no-interference notation alone cannot identify an economy-wide intervention. A difference of mean potential outcomes does not require the cross-world joint distribution. Conditioning on post-treatment migration, survival, enrollment or employment changes the estimand and needs separate justification.

\textbf{Statistical statements.} An estimator, confidence set, forecast or test specifies a sampling law, model class and loss/coverage criterion. Uniform \((1-\alpha)\)-coverage means \(\inf_{m\in\mathcal M}\Pr_m\{\tau(m)\in C_n\}\geq1-\alpha\), or an explicitly asymptotic version. The whole parameter space trivially has coverage, so informativeness requires a stated width, risk or power objective. A forecasting improvement is relative to a named benchmark, information set and evaluation population.

\textbf{Equilibrium and optimization.} An equilibrium comprises feasible plans, optimality, beliefs where needed, budget constraints and all market/resource clearing. The primitives or a stated benchmark define each correspondence \(\mathcal E(p;m)\). Nonemptiness is assumed or proved, never inferred just from compact policy sets. A selected-equilibrium target uses a declared selection rule; a robust target quantifies over all permitted equilibria. Use suprema or infima unless attainment is established. An \(\varepsilon\)-optimal policy is feasible and within \(\varepsilon\) of the finite optimum value. Report an empty feasible set separately.

\textbf{Welfare and accounting.} Normative utility, interpersonal normalization, social weights, numeraire, discounting and the use of public receipts are primitives. Behavioral utility need not coincide with the welfare criterion. Household welfare includes ownership income and rents once; adding profits again to compensating variation generally double-counts them. A monetary equivalent is defined by an explicit transfer and price convention, with monotonicity and range conditions. Average effects, mechanism contributions, transfers and welfare losses are different targets. Infinite sums require convergence and dynamic feasible plans require terminal or no-Ponzi conditions.

\textbf{Source versions.} A working-paper date, revision date and journal publication year are distinguished. A DOI for a journal version is not automatically the DOI of an earlier manuscript. Locators refer to the stated version and printed pages where specified. A metadata-only check does not verify a claim about a full-text conclusion. Every entry's source subsection narrows the provenance claim accordingly.
% END ECONOMICS STATEMENT
\end{document}
